\chapter{The de Branges Sampling Comparison}
\label{v4-arith:w7-c:chapter}

\section{The model space}

Write $F^\sharp(z)=\overline{F(\bar z)}$.
Use the Hermite--Biehler condition $|E^\sharp(z)|<|E(z)|$ for $\Im z>0$.
This convention permits common real zeros of $E$ and $E^\sharp$.
The quotients below use their removable continuations.
For such an $E$, put $\Theta=E^\sharp/E$ and
\[
 H^2=\left\{h\text{ holomorphic on }\mathbb C_+:
      \sup_{y>0}\int_{\mathbb R}|h(x+iy)|^2dx<\infty\right\},
 \qquad K_\Theta=H^2\ominus\Theta H^2.
\]
The de Branges space consists of entire $F$ with $F/E,F^\sharp/E\in H^2$,
and has norm $\|F\|_{H(E)}=\|F/E\|_{L^2}$.
Multiplication by $E$ is the isometry $K_\Theta\to H(E)$.
These conventions are given in \href{https://arxiv.org/html/2301.00421v3}{Suzuki, Sections~2.3--2.4}.

\begin{example}[a zero of arbitrary multiplicity]\label{v4-arith:db:ex:multiple}
Let $A(z)=z^m$ for an integer $m\ge1$. Then
\[
 E=A+iA'=z^{m-1}(z+im),\qquad
 \Theta(z)=\frac{z-im}{z+im},\qquad
 e_0(z)=i\sqrt{\frac m\pi}\frac1{z+im}.
\]
The space $K_\Theta$ is one-dimensional, and $e_0$ is a unit vector because
$\int_{\mathbb R}(x^2+m^2)^{-1}dx=\pi/m$.
Also $e_0(0)=1/\sqrt{\pi m}$ and $\Theta'(0)=-2i/m$.
Thus $\|h\|_{K_\Theta}^2=\pi m|h(0)|^2$ on this space.
The multiplicity enters the sampling measure.
It is compatible with a single normalized basis vector at the distinct zero.
\end{example}

\begin{theorem}[explicit de Branges sampling realization]\label{v4-arith:db:thm:db}
\label{v4-arith:w7-c:thm:Mellin-dB-isometry}
\label{v4-arith:w7-c:cor:ATP-as-dB-positivity}
Assume that all zeros of $\xi$ lie on the critical line.
Put $A(z)=\xi(1/2+iz)$, $E=A+iA'$, and $\Theta=E^\sharp/E$.
Let $\Gamma$ be the distinct real zeros of $A$, with multiplicities $m_\gamma$.
Then $E$ is Hermite--Biehler and
\begin{equation}\label{v4-arith:db:eq:basis}
 e_\gamma(z)=i\sqrt{\frac{m_\gamma}{\pi}}
       \frac{A(z)}{(z-\gamma)E(z)},\qquad \gamma\in\Gamma,
\end{equation}
form an orthonormal basis of $K_\Theta$.
For $f\in\mathscr A=C_c^\infty(\mathbb R)$, with
$h_f(z)=\int f(t)e^{izt}dt$, the series
\begin{equation}\label{v4-arith:db:eq:db-map}
 T f=\sum_{\gamma\in\Gamma}\sqrt{\pi m_\gamma}\,
              h_f(\gamma)e_\gamma\quad\text{in }K_\Theta,
 \qquad \mathcal L f=\pi^{-1/2}E\,Tf\quad\text{in }H(E)
\end{equation}
converges and satisfies
\begin{equation}\label{v4-arith:db:eq:db-pairing}
 \langle\mathcal L f,\mathcal L g\rangle_{H(E)}=q_W(f,g).
\end{equation}
For every $h\in K_\Theta$,
\begin{equation}\label{v4-arith:db:eq:sampling-norm}
 \|Eh\|_{H(E)}^2=\|h\|_{L^2}^2
            =\pi\sum_{\gamma\in\Gamma}m_\gamma|h(\gamma)|^2.
\end{equation}
\end{theorem}
\begin{proof}
The even order-one Hadamard product of $A$ gives, with positive zeros grouped,
\[
 A(z)=A(0)\prod_{\gamma>0}(1-z^2/\gamma^2)^{m_\gamma},\qquad
 L(z):=\frac{A'(z)}{A(z)}
   =\sum_{\gamma>0}m_\gamma
          \left(\frac1{z-\gamma}+\frac1{z+\gamma}\right).
\]
The grouped sum converges locally uniformly off the zeros because
$\sum m_\gamma/(1+\gamma^2)<\infty$.
Evenness removes the linear exponential in the Hadamard product.
There is no zero at the origin: $\zeta(1/2)\ne0$, as follows from
the positive alternating eta series and $1-2^{1/2}<0$.
For $\Im z>0$, each reciprocal has negative imaginary part.
Thus $\Im L(z)<0$ and
\[
 |E(z)|^2-|E^\sharp(z)|^2=-4|A(z)|^2\Im L(z)>0.
\]
This proves the Hermite--Biehler condition.

The boundary norm of $H^2$ uses $du$, so its kernel is
$[2\pi i(\bar w-z)]^{-1}$.
Orthogonal projection onto $K_\Theta$ therefore gives the model kernel
\[
 k_\Theta(w,z)=\frac{1-\overline{\Theta(w)}\Theta(z)}{2\pi i(\bar w-z)}.
\]
Since $\Theta=(1-iL)/(1+iL)$, direct algebra gives
\begin{align*}
 k_\Theta(w,z)
 &=\frac{L(z)-\overline{L(w)}}{
       \pi(\bar w-z)(1+iL(z))(1-i\overline{L(w)})}\\
 &=\sum_{\gamma\in\Gamma}e_\gamma(z)\overline{e_\gamma(w)}.
\end{align*}
The difference of each pair of reciprocal terms cancels its linear denominator.
The resulting series is absolutely locally convergent for $z,w\in\mathbb C_+$.

Near $\gamma$, write $A(z)=(z-\gamma)^{m_\gamma}a(z)$ with $a(\gamma)\ne0$.
Then
\[
 \frac{A(z)}{E(z)}=\frac{z-\gamma}{im_\gamma}+O((z-\gamma)^2),\qquad
 \Theta'(\gamma)=-\frac{2i}{m_\gamma}.
\]
It follows that
$e_\gamma(\gamma)=1/\sqrt{\pi m_\gamma}$ and
$e_\gamma(\delta)=0$ for every distinct zero $\delta$.
The kernel has the finite boundary limit
$k_\Theta(\gamma,z)=e_\gamma(z)/\sqrt{\pi m_\gamma}$.
These boundary evaluations are bounded: the kernels at $\gamma+iy$
converge in norm, as the inner products of these kernels converge to their finite analytic limits.
Consequently the $e_\gamma$ are normalized reproducing kernels and are orthonormal.

For fixed $w\in\mathbb C_+$, the kernel expansion on its diagonal gives
$\sum_\gamma|e_\gamma(w)|^2=k_\Theta(w,w)$.
The orthogonal projection of $k_\Theta(w,\cdot)$ onto their closed span
therefore has the full norm of that kernel.
All reproducing kernels lie in the span, proving completeness.
Parseval now gives \eqref{v4-arith:db:eq:sampling-norm}.
The rapid strip decay of $h_f$ gives
$\sum m_\gamma|h_f(\gamma)|^2<\infty$, so \eqref{v4-arith:db:eq:db-map} converges in norm.
Its coefficients give
$\langle Tf,Tg\rangle=\pi q_W(f,g)$.
Multiplication by $E$ and division by $\sqrt\pi$ prove \eqref{v4-arith:db:eq:db-pairing}.
\end{proof}

The basis and its normalization also appear in
\href{https://arxiv.org/html/2301.00421v3}{Suzuki, equation~(3.5), Proposition~4.1}.
The theorem assumes real zeros before constructing a positive realization.
Conversely, if this $E$ is Hermite--Biehler, a nonreal zero of $A$ is impossible.
At an upper-half-plane zero, $E=iA'$ and $E^\sharp=-iA'$ have equal absolute value.
Real conjugation excludes lower-half-plane zeros as well.
Thus this Hermite--Biehler condition contains the real-zero requirement.

\begin{proposition}[the Gaussian domain obstruction]
\label{v4-arith:db:prop:hardy}
For $f_\lambda(t)=e^{-\lambda t^2/2}$ with $\lambda>0$, one has
$h_{f_\lambda}\notin H^2(\mathbb C_+)$.
Consequently $E h_{f_\lambda}\notin H(E)$.
\end{proposition}
\begin{proof}
The Gaussian Fourier formula gives
\[
 \int_{\mathbb R}|h_{f_\lambda}(x+iy)|^2dx
 =\frac{2\pi^{3/2}}{\sqrt\lambda}e^{y^2/\lambda}.
\]
This is unbounded as $y\to\infty$.
Membership of $E h_{f_\lambda}$ in $H(E)$ would require its quotient
by $E$ to belong to $H^2$.
\end{proof}

For $h\in K_\Theta$, multiplication by $E$ transports the Lebesgue
norm. For a compact arithmetic test, the sampling series $Tf$ is
an explicit map into that model space. The identity
$Tf=h_f$ requires $h_f\in K_\Theta$ and follows from equality of
all samples in the displayed orthonormal basis.
The construction therefore does not identify an arbitrary entire
Fourier transform with a model-space vector.

The signed comparison \eqref{v4-arith:sc:full-comparison} holds
without the real-zero assumption. The positive map $\mathcal L$
requires that assumption before its construction.
A finite-prime Petersson norm on all compact tests is impossible by
Proposition~\ref{v4-arith:sc:two-signs}. On its explicit positive line, the
finite-prime form has the stated $L^2$ isometry.

\begin{proposition}[sampling weights and operator multiplicities]
The sampling model has one basis vector $e_\gamma$ per distinct
zero $\gamma$. The diagonal operator $e_\gamma\mapsto\gamma e_\gamma$
on its maximal domain therefore has simple eigenvalues.
The weight $m_\gamma$ in the sampling identity does not give an
$m_\gamma$-dimensional eigenspace.
The occurrence carrier $K_Z$ does have the required multiplicities
for the determinant \eqref{v4-arith:sc:normalized-det-formula}.
\end{proposition}
\begin{proof}
The complete orthonormal basis determines every eigenspace of the
first diagonal operator. Its coordinates are indexed by distinct
zeros, so each eigenspace is one-dimensional.
The coordinates of $K_Z$ instead contain each zero occurrence.
Changing a scalar product weight changes the norm of a vector,
but not the dimension of its eigenspace.
\end{proof}
